\section{Response to Reviewer \texttt{hBhv}}

\subsection*{Summary}

A placebo-anchored transport framework is proposed to address the challenges in representing broader patient populations of randomized clinical trials.


\subsection*{Summary of Strengths}

\begin{itemize}
    \item The research problem and its significance are clearly articulated.
    \item An experiment is conducted to analyze the strength of the proposed methods, and the analysis is clearly presented.
\end{itemize}


\subsection*{Summary of Weaknesses}


\subsubsection{Motivation relative to existing literature.}  
A literature review is presented and the weaknesses of existing methods are discussed. However, the motivation for the proposed method is not clearly articulated in terms of how it directly addresses the identified shortcomings.

{\color{blue}
We thank the reviewer for this feedback. In the revised manuscript, we have strengthened the connection between the identified shortcomings of existing methods and the design of our approach. Specifically, we now explicitly articulate that:
\begin{itemize}
\item Standard transport methods require conditional exchangeability $(Y(a) \perp C \mid X)$, which is often violated in multi-site RCTs due to unmeasured site-level confounding;
\item Meta-analytic approaches assume exchangeability across trials, which may not hold under systematic covariate shift;
\item Our placebo-anchored framework directly addresses these limitations by (i) using target placebo outcomes to calibrate baseline risk without requiring treated target data, and (ii) imposing a structured regularity condition (Assumption A5) rather than full exchangeability.
\end{itemize}
The introduction has been revised to make this motivation explicit.
}

\subsubsection{Lack of benchmarking against state-of-the-art methods.}  
The experiments only present comparisons among different settings of the proposed method. There is no benchmarking against state-of-the-art methods in the literature, making it unclear whether the proposed approach outperforms existing techniques, despite the analysis being interesting.

{\color{blue}
We thank the reviewer for this important comment and agree that benchmarking against established methods in the literature is necessary to contextualize the proposed approach. In response, we have substantially expanded the experimental section to include comparisons with 10 methods spanning ablation baselines, transport estimators, and oracle references.

Specifically, we now benchmark against:
\begin{itemize}
\item \textbf{Ablation baselines}: ProxyOnly (no target data), AnchorOnly (target data only, no proxy borrowing), and AnchorPlugin (outcome-model transport with baseline anchoring);
\item \textbf{Reweighting-based transport}: IPWTransport (inverse probability weighting for site membership);
\item \textbf{Outcome-regression transport}: OM-Transport, which pools source trials and predicts on target covariates;
\item \textbf{Moment-based transport}: EntropyBalancing, which reweights sources to match target covariate moments;
\item \textbf{Proposed methods}: Proposed (Option A with target treated data), Proposed-CF (cross-fitted variant), and Proposed-B (Option B for disconnected regime with $m_1 = 0$).
\end{itemize}

All methods are evaluated over 100 Monte Carlo replications across three comprehensive sweeps: (i) target budget vs.\ dimensionality ($p \in \{10, 20, 50, 100\}$), (ii) source site scaling ($C \in \{2, 5, 10, 20, 50\}$), and (iii) Assumption A5 violation sensitivity.

\medskip
\noindent
\textbf{Alignment with intended data regimes.}
Methods have heterogeneous data requirements. To ensure fair comparison, we distinguish:
\begin{itemize}
\item the \emph{disconnected/restricted regime} ($m_1=0$), where only target placebo is available---transport baselines (IPWTransport, OM-Transport, EntropyBal), AnchorPlugin, ProxyOnly, and Proposed-B are evaluated;
\item the \emph{oracle regime} ($m_1>0$), where target treated outcomes exist---all 10 methods including TargetOnlyDR, Proposed, and Proposed-CF are compared.
\end{itemize}

\medskip
\noindent
\textbf{Key findings.}
In the \textbf{disconnected regime} ($m_1=0$), Proposed-B matches or outperforms transport baselines across dimensions and source site counts, achieving the best PEHE in most conditions.
In the \textbf{oracle regime} ($m_1>0$), Proposed dramatically outperforms all baselines: at $m_1=50$, it achieves PEHE of $0.67 \pm 0.17$ (vs.\ $2.11$ for OM-Transport, $2.82$ for AnchorOnly), Spearman correlation of $0.98 \pm 0.01$ (vs.\ $0.83$ for transport baselines), and policy regret of $0.02 \pm 0.01$ (vs.\ $0.28$ for transport baselines).

These results demonstrate that the proposed method provides substantial gains over state-of-the-art transport approaches when target treated data is available, and remains competitive when operating in the more challenging disconnected regime.
}



\subsubsection{Missing conclusion section.}  
A conclusion section is missing that would highlight the main contributions and findings of the proposed work.

{\color{blue}
We thank the reviewer for this comment. We apologize for the confusion caused by the original typesetting. In the revision, we have added a dedicated \emph{Conclusion} section and improved the section formatting (clear heading placement and cross-referencing) so it is easy to locate. The new conclusion succinctly summarizes our main contributions---(i) the placebo-anchored proxy--gold framing for multi-trial transport under covariate shift, (ii) the connected-target DR estimator (\texttt{Proposed-CF}) and the disconnected-target screen--then--transport procedure (\texttt{Proposed-B}) with the updated Assumption~A6, and (iii) the key empirical findings on calibration, CATE accuracy, and graceful degradation under assumption violations. We also include a brief discussion of limitations and future directions (diagnostics and sensitivity analysis for screening-valid transportability, extensions to richer endpoints/learners, and uncertainty quantification in small-gold regimes).
}


\clearpage